\chapter{Implementation}
\label{chap:implementation}

The software architecture developed for this thesis realizes the integration of Federated Dynamic Averaging with parameter-efficient fine-tuning via FDA-RoLoRA, utilizing an internal dual-track threshold mechanism to manage coordinate subspace variance, open-sourced under the \textbf{FDA-RoLoRA} repository. This chapter details the software architecture, the strategies for data partitioning across clients, the configuration of the targeted foundation models, the mechanics of the training pipeline, and the high-performance computing environment utilized for experimental evaluation.

\section{Software Architecture}
\label{sec:impl_architecture}

The framework is developed in Python, utilizing the PyTorch deep learning framework as the foundational library for automatic differentiation and tensor operations. The codebase is built as a direct extension of the FDA-Opt repository \cite{theologitis2025fda},\footnote{\url{https://github.com/michaeltheologitis/FDA-Opt}} inheriting its robust core federated orchestration infrastructure. To implement the adaptation strategies, the system heavily relies on the Hugging Face Transformers library in conjunction with the Parameter-Efficient Fine-Tuning (PEFT) library, which has established itself as the research standard for LoRA-based methods. The complete executable software architecture developed for this thesis has been open-sourced as \textbf{FDA-RoLoRA}.\footnote{\url{https://github.com/Asterinos1/FDA-RoLoRA}}

A key practical constraint faced during this research was the absence of a publicly available code repository for RoLoRA \cite{chen2025robust}, which only noted using the FederatedScope-LLM platform \cite{kuang2023federatedscope}. Given that we had direct access to the FDA-Opt codebase \cite{theologitis2025fda}, the design choice was between porting FDA's dynamic variance monitoring and threshold evolution logic into the heavy, event-driven architecture of FederatedScope-LLM, or implementing RoLoRA's alternating optimization protocol within the existing FDA-Opt framework. We chose the latter approach, as implementing RoLoRA inside the lightweight FDA-Opt core was substantially more tractable and transparent. This decision afforded complete architectural visibility and direct control over low-rank tensor serialization and threshold updates, without the abstraction overhead of a third-party framework.


To manage the complexities of federated learning, the system follows a clear three-tier architecture comprising a Server Orchestrator, Client Nodes, and a Communication Interface. The Server Orchestrator assumes a centralized role; it maintains the global LoRA parameter matrices $A$ and $B$, tracks the current threshold states $\theta_A$ and $\theta_B$ for synchronization, manages server-side optimizer momentum buffers, and increments a global round counter responsible for the RoLoRA alternating optimization toggle. The Client Nodes represent the edge devices or data silos. During each communication round, a subset of sampled clients receives the active parameter block ($A$ or $B$), performs local SGD on their private datasets across dynamic local epoch intervals, and returns the accumulated parameter drift vectors ($\Delta_i$) to the server, which centrally evaluates cohort variance and orchestrates the aggregation step. Finally, the Communication Interface orchestrates the logical data exchange between clients and the server. In our simulation testbed, client-server communication is modeled in-silico: rather than incurring physical network socket overhead, parameter drift vectors ($\Delta_i$) are calculated via direct in-memory tensor subtractions against the round baseline snapshot. Communication bandwidth consumption is evaluated analytically by measuring the exact serialized byte footprint of the active low-rank coordinate tensors transmitted per client per round (e.g., $4.69\,\text{MB}$ for rank $r=4$ adapters in single-precision float32). Unlike conventional federated learning frameworks that transmit full model weights, this design yields a massive reduction in bandwidth requirements. 

Crucially, the alternating optimization protocol—the toggling mechanism inherent to the RoLoRA methodology—is strictly managed server-side. This centralization guarantees global consistency across all participating clients, ensuring that during any given round, every active client is optimizing and communicating updates for the same LoRA matrix block.


\section{Data Partitioning}
\label{sec:impl_data_partitioning}

The software architecture implements client data partitioning following the Dirichlet allocation protocol defined in Section~\ref{sec:dirichlet}. During environment initialization, the partitioning module loads the targeted GLUE corpus. To address the significant disparity in scale across benchmark tasks (ranging from under $70{,}000$ training examples in SST-2 to nearly $400{,}000$ in MNLI) and mitigate prohibitive simulation runtimes on the shared GPU cluster, the framework adopts an empirical dataset ceiling. Specifically, the training volume is capped at an upper ceiling of $25{,}000$ examples via deterministic pseudo-random subsampling, while smaller corpora and all validation evaluation sets are preserved in full. Preliminary pilot investigations confirmed that this sample budget provides sufficient linguistic diversity and gradient support for adapter fine-tuning without compromising convergence stability or task performance, while keeping multi-seed simulation horizons computationally tractable. Following this volume standardization, the module segments the active corpus into $K$ isolated local partitions according to the configured Dirichlet concentration parameter $\alpha$. Each simulated client receives a dedicated local partition and batch iterator, enforcing strict data isolation during local training without cross-client leakage.

\section{Simulation Controls and Early Stopping}
\label{sec:impl_early_stopping}

Beyond dataset scale, substantial wall-clock latency accumulates during simulation because client cohorts ($C \in \{20, 50\}$) are evaluated sequentially on the GPU to respect accelerator memory limits. To prevent excessive simulation runtimes across multi-seed sweeps and avoid computing redundant rounds on converged models, the framework incorporates an \textbf{accuracy-gated early stopping controller} into the Server Orchestrator. At the conclusion of each communication round, the server evaluates the consensus model on the centralized task validation split. The controller governs early termination via four complementary mechanisms:

\begin{enumerate}
    \item \textbf{Trailing Sliding-Window Smoothing ($W$)}: In federated optimization under statistical heterogeneity, validation trajectories exhibit stochastic oscillations due to client sampling and non-IID drift. To prevent premature halts triggered by transient dips, the controller tracks a trailing moving average of validation accuracy across a sliding window of $W$ rounds.
    \item \textbf{Accuracy Activation Gate ($A_{\text{gate}}$)}: Because alternating coordinate optimization requires several initial rounds to calibrate both factor matrices ($A$ and $B$), early stopping is strictly disarmed during early training. The patience counter activates only after the smoothed validation accuracy crosses a predefined milestone threshold ($A_{\text{gate}}$).
    \item \textbf{Patience and Minimum Improvement ($\Delta_{\min}, P$)}: Once the activation gate is breached, the controller monitors validation progress over a patience window of $P$ consecutive rounds. If the moving-average accuracy fails to improve by at least $\Delta_{\min} = 0.5\%$, the training loop terminates cleanly ahead of $R_{\max}$.
    \item \textbf{Mandatory Round Guard ($R_{\min}$)}: To ensure that models explore sufficient alternating coordinate cycles regardless of metric volatility, the engine enforces a strict round floor ($R_{\min}$) before any termination signal can trigger.
\end{enumerate}

This controller functions as an external simulation management harness, accelerating cluster turnaround times without modifying the underlying federated optimization dynamics. The specific hyperparameter configurations deployed across each GLUE benchmark task and partitioning regime are detailed in Section~\ref{sec:eval_early_stopping}.

\section{Model Configuration}
\label{sec:impl_model_config}

The experimental implementation targets transformer-based foundation models. In this work, the empirical evaluation focuses specifically on \texttt{roberta-base} \cite{liu2019roberta} (comprising approximately 125 million parameters), evaluated on sequence classification tasks from the GLUE benchmark \cite{wang2019glue}. While the original RoLoRA paper \cite{chen2025robust} evaluated primarily on \texttt{roberta-large} ($355\text{M}$ parameters), we deliberately chose \texttt{roberta-base} to accommodate our institutional hardware budget ($2 \times$ NVIDIA A10 GPUs) and project turnaround time limits. This choice cuts VRAM allocation and backward pass latency by more than half, enabling thorough multi-seed sweeps across 20--50 clients while maintaining representative transformer attention geometry. For \texttt{roberta-base}, the low-rank adapters are systematically injected into the query and value projection matrices ($W_q$ and $W_v$) within each self-attention layer.

Across the model configuration, the LoRA rank is strictly constrained to $r=4$ (and $r=8$), accompanied by scaling factor $\alpha_{\text{lora}} = 8$ (and $\alpha_{\text{lora}} = 16$) to maintain activation magnitudes. During the federated fine-tuning process, all original pretrained foundation model weights as well as the downstream sequence classification head are kept completely frozen. The only trainable parameters across all communication rounds are the low-rank adapter matrices $A$ and $B$. This setup provides a profound advantage in communication cost, as communicating only the compact low-rank drift matrices rather than the full 125M parameter set yields orders-of-magnitude reductions in client-server bandwidth requirements. While broader foundation models (such as RoBERTa-Large or autoregressive decoders like Llama-2-7B) represent natural targets for future scaled deployment, our implementation and empirical study focus on \texttt{roberta-base} across GLUE tasks.

\section{Training Pipeline}
\label{sec:impl_training_pipeline}

The operational training pipeline translates the theoretical FDA-RoLoRA protocol (formulated in Algorithm~\ref{alg:fda_rolora}) into an efficient, stateful PyTorch execution loop. At the start of communication round $t$, the Server Orchestrator determines the active parameter track by evaluating round parity ($t \pmod 2$), activating track $B$ on even rounds and track $A$ on odd rounds. To enforce alternating coordinate isolation in software, the orchestrator traverses the model parameter tree (\texttt{model.named\_parameters()}) and dynamically toggles the PyTorch autograd flags: parameters matching the active adapter key are assigned \texttt{requires\_grad = True}, while the dormant adapter block, the downstream classification head, and all backbone foundation model parameters are assigned \texttt{requires\_grad = False}. The server snapshots the active parameters ($W_{\text{act}}^{(0)}$) as the baseline for tracking parametric drift, binds the corresponding independent server optimizer instance, loads the active threshold state ($\theta_B$ or $\theta_A$), and synchronizes the active weights to the sampled client cohort $S_t \subset \{1, \dots, K\}$. Maintaining two persistent \texttt{torch.optim} instances across rounds keeps optimizer momentum states strictly partitioned, preventing cross-phase state contamination.

Within each simulated communication round, sampled clients are evaluated sequentially on the allocated GPU accelerator, executing local SGD over their private data partitions using a configurable local learning rate $\eta_c$ and local batch size $B$. This sequential simulation design faithfully models the logical decentralized execution of independent client nodes while strictly constraining GPU VRAM footprint to a single active client instance at any given instant. Because inactive weights have autograd disabled, PyTorch allocates gradient and optimizer memory exclusively for the active low-rank coordinates, minimizing client VRAM consumption. Local training proceeds across discrete epoch boundaries bounded by runtime-configurable limits $[e_{\min}, e_{\max}]$ (typically configured within $[1, 20]$ across our empirical foundation model benchmarks). For each epoch $e \ge e_{\min}$, the framework extracts each client's local parameter drift $\Delta_i^{(e)} = W_{i,e} - W_{\text{act}}^{(0)}$ via direct in-memory tensor subtraction against the round baseline. The orchestrator flattens each drift tensor, stacks them via \texttt{torch.stack}, and computes the exact subspace variance $\text{Var}_{\text{subspace}}(e)$ across the cohort following Eq.~\ref{eq:subspace_var} in single-digit milliseconds using vectorized PyTorch operations (\texttt{torch.norm}). The loop automatically halts at $e^* = e_{\min}$ during the first round of each matrix track to record baseline variance and initialize the threshold via Eq.~\ref{eq:bootstrap_init}. In subsequent rounds, the server dynamically halts local training at the earliest epoch $e^*$ where $\text{Var}_{\text{subspace}}(e) > \theta_{\text{current}}$, or upon reaching the ceiling $e^* = e_{\max}$ (Eq.~\ref{eq:halt_cond}).

Following inner loop termination, the server aggregates the gathered drifts into the pseudo-gradient $G_t = -\frac{1}{|S_t|} \sum_{i \in S_t} \Delta_i^{(e^*)}$ (Eq.~\ref{eq:pseudo_grad}). Rather than performing manual weight substitution, the server restores active parameters to $W_{\text{act}}^{(0)}$, injects $G_t$ directly into the parameter \texttt{.grad} buffers, and invokes \texttt{optimizer.step()} using the active server optimizer instance. With standard SGD ($\eta_s = 1.0$), this operation executes coordinate-wise Federated Averaging (Eq.~\ref{eq:server_fedavg}) while preserving a decoupled interface compatible with adaptive server optimizers (such as AdamW, Eq.~\ref{eq:adamw}). The server then recalibrates the active track's synchronization threshold using the Standard midpoint scaling rule (Eq.~\ref{eq:std_update}) while keeping the dormant threshold frozen in memory.

To prevent accelerator memory fragmentation across hundreds of simulation rounds on the A10 GPUs, the pipeline enforces strict memory hygiene: gradient buffers are deallocated via \texttt{optimizer.zero\_grad(set\_to\_none=True)}, intermediate computation graphs and drift tensors are explicitly detached, and the orchestrator periodically flushes the CUDA cache using \texttt{torch.cuda.empty\_cache()} before incrementing the round counter ($t \leftarrow t + 1$) to launch the complementary matrix phase.

\section{Deployment Environment}
\label{sec:impl_deployment}

The development and execution environment for the FDA-RoLoRA framework is designed to ensure strict software reproducibility and computational efficiency. The software stack is standardized across all nodes,\footnote{All experiments were conducted using Python 3.10.12, PyTorch 2.5.1+cu121, Hugging Face Transformers 5.5.4, and PEFT 0.19.1, running CUDA 12.1 with cuDNN 9.1 on Ubuntu Linux.} pinning core libraries to guarantee consistent autograd graph mechanics, deterministic low-rank tensor serialization, and reproducible pseudo-gradient aggregation across simulation runs.
 
Production experiments are executed on the Polytechnix server of the SoftNet Lab at the Technical University of Crete. The server is equipped with enterprise-grade NVIDIA accelerators ($2 \times$ NVIDIA Ampere A10 GPUs with $24\,\text{GB}$ VRAM each), $2 \times$ Intel Xeon Silver 4310 CPUs ($24$ cores / $48$ threads total at $2.10\,\text{GHz}$), and $256\,\text{GB}$ of RAM, providing the necessary computational throughput for fine-tuning foundation models.
 
A suite of custom orchestration bash and Python scripts manages the execution lifecycle. Rather than spawning parallel client processes within a single simulation (which would saturate GPU VRAM), multi-process parallelism is leveraged at the cluster level: the orchestration scripts dispatch distinct experimental configurations (spanning multiple random seeds, baseline comparisons, and Dirichlet partitioning regimes) concurrently across the two available NVIDIA A10 GPUs. Within each running simulation job, the script coordinates dynamic synchronization rounds, logs tracking metrics, and automates the extensive parameter sweeps required for empirical evaluation and ablation analysis.
